% Separate book (split from the main book 001_foundations.tex on 2026-11-28).
% Build from this folder: pdflatex category_theory_topics.tex ; it reuses ../000_preamble.tex via ../BEGIN.tex.
\def\mypath{../}
\def\topic{Topics in category theory}
\ifdefined\MyHeaderIncluded
\else
\def\qwertyflag{true}

\ifdefined\MyHeaderIncluded

\else
  \newcommand{\foobar}{FooBar}
  \def\MyHeaderIncluded{true}
% if is till the end of the file

\documentclass[10pt,a4paper]{article}
\usepackage[utf8]{inputenc}
%\usepackage[german]{babel}
%\usepackage[cyrillic]{babel}
\usepackage[T1]{fontenc}
%\usepackage[T2A]{fontenc}
%\usepackage[english, russian]{babel}
\usepackage{amsmath}
\usepackage{amsfonts}
\usepackage{amssymb}
\usepackage{amsthm}
\usepackage{mathtools} %coloneqq

\usepackage{bussproofs} %

%\renewcommand{\thesubsection}{\arabic{subsection}}
%\newtheorem{thm}{Theorem}[section]
%\renewcommand{\thethm}{\arabic{subsection}.\arabic{thm}}

\usepackage{thmtools}
\usepackage{hyperref}
\hypersetup{
    colorlinks=true,
    linkcolor=blue,
    filecolor=magenta,      
    urlcolor=cyan,
    pdftitle={Overleaf Example},
    %pdfpagemode=FullScreen,
    }
%\usepackage{stix-type1}%\turnediota
\theoremstyle{definition}

\newtheorem{constructor}{Constructor}[section]
\newtheorem{iaxiom}{Axiom}
\newtheorem{axiom}{Axiom}
\newtheorem*{axiom*}{Axiom}
\newtheorem{definition}{Definition}[section]
\newtheorem*{definition*}{Definition}
\newtheorem{remark}{Remark}[section]
\newtheorem*{remark*}{Remark}
\newtheorem{recall}{Recall}[section]
\newtheorem*{recall*}{Recall}

\newtheorem{lemma}{Lemma}[section]
\newtheorem{theorem}{Theorem}[section]
\newtheorem*{theorem*}{Theorem}
\newtheorem{note}{Note}[section]
\usepackage{etoolbox}

%\newcommand*{\@Dashv}[2]{%
%  \reflectbox{$\m@th#1#2$}%
%}
\newcommand*{\separation}[3]{% 1)binder, 2)set, 3)condition
  \{#1\in #2: #3\}
}

% https://www.freecodecamp.org/news/simple-regex-tricks-for-beginners-3acb3fa257cb/
%  \\mathcal\{T\} |---> \topology
%  \\mathcal\{([^\}]*)\}   |--->  $1
\newcommand{\AND}{\ \ \&\ \ }
\newcommand{\any}{\mathrm{\ any}}
\newcommand{\borel}{\mathcal{B}}
\newcommand{\can}{\mathrm{can}}
\newcommand{\canonize}{\mathrm{canonize}}
\newcommand{\closure}{\mbox{cl}}
\newcommand{\complex}{\mathbb{C}}
\newcommand{\converges}{\longrightarrow}
\newcommand{\countable}{\mbox{\ countable}}
\newcommand{\ldefined}{\mbox{\ defined}}
\newcommand{\dom}{\mbox{dom}}
\newcommand{\Dom}{\mbox{Dom}}
\newcommand{\Done}{\noindent \textbf{Done.}\\}
\newcommand{\lexists}{\mbox{\ exists}}
\newcommand{\lforall}{\mbox{\ forall}}
\newcommand{\lunique}{\mbox{\ unique}}
\newcommand{\family}{\mathrm{\ family}}
\newcommand{\Fm}{\mathrm{Fm}}
\newcommand{\FuncSymb}{\mathrm{FuncSymb}}
\newcommand{\function}{\mathrm{\ function}}
\newcommand{\functional}{\mbox{\ functional}}
\newcommand{\IFF}{\ \Longleftrightarrow\ }
\renewcommand{\iff}{\ \longleftrightarrow\ }
\newcommand{\IMPLIES}{\ \Longrightarrow\ }
\renewcommand{\implies}{\ \longrightarrow\ }
\newcommand{\Inf}{\mathrm{Inf}}
\newcommand{\injective}{\mbox{\ injective}}
\newcommand{\injection}{\mbox{\ injection}}
\newcommand{\interior}{\mbox{int}}
\newcommand{\intersection}{\bigcap}
\newcommand{\isempty}{\mathrm{\ empty}}
\newcommand{\LB}{\mathrm{LB}}
\newcommand{\nghbset}{\mathcal{N}}
\newcommand{\nonempty}{\mathrm{\ nonempty}}
\newcommand{\Ord}{\mathrm{Ord}}
\newcommand{\power}{\mathcal{P}}
\newcommand{\pr}{\mathrm{pr}}
\newcommand{\fst}{\mathrm{fst}}
\newcommand{\snd}{\mathrm{snd}}
\newcommand{\PredSymb}{\mathrm{PredSymb}}
\newcommand{\Proof}{\noindent \textbf{Proof.}\\}
\newcommand{\ran}{\mathrm{ran}}
\newcommand{\Ran}{\mathrm{Ran}}
\newcommand{\real}{\mathbb{R}}
\newcommand{\relation}{\mathrm{\ relation}}
\newcommand{\set}{\mathrm{\ set}}
\newcommand{\singleton}{\mathrm{\ singleton}}
\newcommand{\semantics}{\mathrm{semantics}}
\newcommand{\semialgebra}{\mathcal{S}}
\newcommand{\Sup}{\mathrm{Sup}}
\newcommand{\surjective}{\mbox{\ surjective\ on\ }}
\newcommand{\Term}{\mathrm{Term}}
\newcommand{\Type}{\mathrm{Type}}
\newcommand{\topology}{\mathcal{T}}
\newcommand{\total}{\mbox{\ total\ on\ }}
\newcommand{\transitive}{\mathrm{\ transitive}}
\newcommand{\union}{\bigcup}
\newcommand{\lif}{\mathrm{if}}
\newcommand{\lthen}{\mathrm{then}}
\newcommand{\lelse}{\mathrm{else}}
\newcommand{\lexception}{\mathrm{exception}}
\newcommand{\lundef}{\mathrm{undef}}
\newcommand{\univocal}{\mathrm{\ univocal}}
\newcommand{\univocalat}{\mathrm{\ univocal\ at\ }}
\newcommand{\UB}{\mathrm{UB}}
\newcommand{\Var}{\mathrm{Var}}

\newcommand{\notvarphi}{\phi}
\renewcommand{\phi}{\varphi}

\newcommand{\comment}[1]{\quad\mbox{(#1)}}
% local commands
%\preto\section{\undef\blipo} %etoolbox
\preto\subsection{\undef\blipo} %etoolbox
% https://tex.stackexchange.com/questions/105684/limit-macro-newcommand-scope-to-section
\usepackage{graphicx}
\usepackage{authblk}

\setlength\parindent{0pt} % whole doc of \noindent

%\usepackage{mathgl}

\fi % guard
\ifdefined\topic
\else
\def\topic{No Titel}
\fi
\title{\topic}
\author[1]{Georgy Dunaev}
\begin{document}
\maketitle
\fi
\section{Topics in category theory}
\begin{note}[Origin of this text]
This text was drafted by an automated assistant (2026-11-15) and has not been reviewed by the author.
It is not formalised in the Obelisk proof assistant. Treat it as a reading list with summaries, not as
proved results.
\end{note}

\subsection{Locally Small Categories, the Yoneda Lemma, and Pointwise Kan Extensions}
\label{subsec:yoneda_presheaves_and_kan_extensions}

A locally small category $\mathcal{C}$ with a proper class of objects $\operatorname{Ob}(\mathcal{C}) \subseteq \V$ represents each hom-set $\operatorname{Hom}_{\mathcal{C}}(X, Y) \in \V$ via ordered pairs $\pair{\pair{X}{Y}}{f}$. For each object $A \in \operatorname{Ob}(\mathcal{C})$, the representable presheaf $h_A(X) = \operatorname{Hom}_{\mathcal{C}}(X, A)$ acts contravariantly by precomposition. Given any set-valued presheaf $F\colon \mathcal{C}^{\mathrm{op}} \to \mathbf{Set}$ encoded at the conglomerate level, the Yoneda evaluation map
\[
\Phi_{A,F}\colon \operatorname{Nat}(h_A, F) \longrightarrow F(A), \qquad \Phi_{A,F}(\eta) = \eta_A(1_A),
\]
is a bijection with inverse $\Psi_{A,F}(u)_X(f) = F(f)(u)$ for all $f \in \operatorname{Hom}_{\mathcal{C}}(X, A)$. Specializing to $F = h_B$ shows that the Yoneda embedding $Y(A) = h_A$ is full and faithful: $\operatorname{Nat}(h_A, h_B) \cong \operatorname{Hom}_{\mathcal{C}}(A, B)$. Furthermore, along a functor $K\colon \mathcal{C} \to \mathcal{D}$ and diagram $F\colon \mathcal{C} \to \mathcal{E}$ with graphs $\pair{c}{K(c)}$ and $\pair{c}{F(c)}$, the pointwise left and right Kan extensions are given by the comma-category colimit and limit formulas
\[
(\operatorname{Lan}_K F)(d) \cong \varinjlim_{(K \downarrow d)} F \circ \pi, \qquad
(\operatorname{Ran}_K F)(d) \cong \varprojlim_{(d \downarrow K)} F \circ \pi,
\]
satisfying the universal adjunction identities $\operatorname{Nat}(\operatorname{Lan}_K F, G) \cong \operatorname{Nat}(F, G \circ K)$ and $\operatorname{Nat}(G \circ K, F) \cong \operatorname{Nat}(G, \operatorname{Ran}_K F)$.

\subsection{Elementary Topoi, Subobject Classifiers, and Lawvere-Tierney Topologies}
\label{subsec:topos_subobject_classifier_lawvere_tierney}

An elementary topos $\mathcal{E}$ is a finitely complete, cartesian closed category equipped with a subobject classifier $\top\colon 1 \to \Omega$ such that for every monomorphism $m\colon S \rightarrowtail X$ there exists a unique characteristic morphism $\chi_m\colon X \to \Omega$ making the diagram
\[
\begin{aligned}
S &\longrightarrow 1, \qquad m\colon S \rightarrowtail X, \qquad \chi_m \circ m = \top \circ {!_S}, \\
S &\cong \cc{\pair{x}{u} \in X \times 1}{\chi_m(x) = \top(u)}
\end{aligned}
\]
a pullback square over ordered pairs $\pair{x}{u}$. A \emph{Lawvere-Tierney modal topology} on $\mathcal{E}$ is an endomorphism $j\colon \Omega \to \Omega$ satisfying three equational axioms:
\begin{enumerate}
  \item \textbf{Truth Preservation}: $j \circ \top = \top$.
  \item \textbf{Idempotence}: $j \circ j = j$.
  \item \textbf{Meet Preservation}: $j \circ \wedge = \wedge \circ (j \times j)$.
\end{enumerate}
Each Lawvere-Tierney topology $j$ induces a closure operator $S \mapsto \overline{S}^j$ on $\operatorname{Sub}_{\mathcal{E}}(X)$ with characteristic morphism $j \circ \chi_m$, and a subobject $m\colon S \rightarrowtail X$ is $j$-closed precisely when $j \circ \chi_m = \chi_m$.

\subsection{Monads, Eilenberg-Moore and Kleisli Categories, and Beck Monadicity}
\label{subsec:monads_eilenberg_moore_kleisli_beck}

A monad on a category $\mathcal{C}$ is a triad $\pair{T}{\pair{\eta}{\mu}}$ encoded via ordered pairs, where $T\colon \mathcal{C} \to \mathcal{C}$ is an endofunctor and $\eta\colon 1_{\mathcal{C}} \Rightarrow T$ (unit) and $\mu\colon T^2 \Rightarrow T$ (multiplication) are natural transformations satisfying the unit and associativity coherence identities:
\[
\mu \circ T\eta = 1_T = \mu \circ \eta T, \qquad \mu \circ T\mu = \mu \circ \mu T.
\]
Each monad $(T, \eta, \mu)$ canonically determines two extremal adjunctions in the category of adjunctions generating $T$:
\begin{enumerate}
  \item \textbf{Eilenberg-Moore Category of $T$-Algebras ($\mathcal{C}^T$, Terminal Resolution)}: Objects are pairs $\pair{A}{h}$ where $h\colon T(A) \to A$ satisfies $h \circ \eta_A = 1_A$ and $h \circ T(h) = h \circ \mu_A$; morphisms $f\colon \pair{A}{h} \to \pair{B}{k}$ satisfy $f \circ h = k \circ T(f)$.
  \item \textbf{Kleisli Category of Free $T$-Algebras ($\mathcal{C}_T$, Initial Resolution)}: Objects are the objects of $\mathcal{C}$, morphisms $f\colon X \rightsquigarrow Y$ in $\mathcal{C}_T$ are morphisms $f\colon X \to T(Y)$ in $\mathcal{C}$ encoded as $\pair{\pair{X}{Y}}{f}$, and Kleisli composition is given by $g \odot f = \mu_Z \circ T(g) \circ f$ with identity $\eta_X$.
\end{enumerate}
Given any adjunction $F \dashv U\colon \mathcal{D} \to \mathcal{C}$ inducing $T = UF$, the canonical comparison functor $K\colon \mathcal{D} \to \mathcal{C}^T$ sends $D \mapsto \pair{U(D)}{U(\varepsilon_D)}$. By Beck's Monadicity Theorem, $K$ is an equivalence of categories if and only if $U$ is conservative (reflects isomorphisms) and $\mathcal{D}$ has and $U$ preserves coequalizers of all $U$-split parallel pairs.

\subsection{Cartesian Closed Categories, Exponential Objects, and Curry-Howard-Lambek Equivalence}
\label{subsec:ccc_exponentials_curry_howard_lambek}

A category $\mathcal{C}$ with finite products (terminal object $1$ and binary product $A \times B$ encoded via ordered pairs $\pair{a}{b}$) is \emph{Cartesian closed} if for every pair of objects $A, B \in \operatorname{Ob}(\mathcal{C})$ there exists an \emph{exponential object} $B^A$ together with a canonical \emph{evaluation morphism} $\operatorname{ev}_{A,B}\colon B^A \times A \to B$ satisfying the universal property of exponentiation: for every morphism $g\colon C \times A \to B$, there exists a unique \emph{exponential transpose} (curried morphism) $\lambda(g)\colon C \to B^A$ such that
\[
\operatorname{ev}_{A,B} \circ (\lambda(g) \times 1_A) = g.
\]
Conversely, given $h\colon C \to B^A$, its \emph{uncurried morphism} is $\operatorname{uncurry}(h) = \operatorname{ev}_{A,B} \circ (h \times 1_A)\colon C \times A \to B$, establishing a natural bijection
\[
\operatorname{Hom}_{\mathcal{C}}(C \times A, B) \cong \operatorname{Hom}_{\mathcal{C}}(C, B^A)
\]
whose two-sided inverse equations mirror the $\beta$- and $\eta$-conversion laws of the simply typed $\lambda$-calculus:
\begin{enumerate}
  \item \textbf{Categorical $\beta$-Reduction Law}: $\operatorname{uncurry}(\lambda(g)) = \operatorname{ev}_{A,B} \circ (\lambda(g) \times 1_A) = g$.
  \item \textbf{Categorical $\eta$-Expansion Law}: $\lambda(\operatorname{uncurry}(h)) = \lambda(\operatorname{ev}_{A,B} \circ (h \times 1_A)) = h$.
\end{enumerate}
Under the Curry-Howard-Lambek correspondence, Cartesian closed categories provide the sound and complete categorical semantics for intuitionistic propositional logic and the simply typed $\lambda$-calculus, interpreting contexts $\Gamma$ as finite products, implication $A \supset B$ as the exponential object $B^A$, modus ponens as $\operatorname{ev}_{A,B}$, and deduction-theorem abstraction as currying $\lambda(-)$.

\subsection{Grothendieck Fibrations, Cartesian Morphisms, Cleavages, and the Grothendieck Construction}
\label{subsec:grothendieck_fibrations_cleavages_construction}

Let $p\colon \mathcal{E} \to \mathcal{B}$ be a functor between categories whose object and morphism graphs are encoded via ordered pairs. For each base object $I \in \operatorname{Ob}(\mathcal{B})$, the \emph{fiber category} $\mathcal{E}_I = p^{-1}(I)$ consists of objects $X \in \operatorname{Ob}(\mathcal{E})$ with $p(X) = I$ and vertical morphisms $v\colon X \to X'$ satisfying $p(v) = 1_I$. A morphism $\varphi\colon X \to Y$ in $\mathcal{E}$ lying over $u = p(\varphi)\colon I \to J$ in $\mathcal{B}$ is \emph{$p$-Cartesian} if for every object $Z \in \operatorname{Ob}(\mathcal{E})$, every morphism $\psi\colon Z \to Y$ in $\mathcal{E}$, and every base factorization $p(\psi) = u \circ w$ with $w\colon p(Z) \to I$, there exists a unique mediating morphism $\theta\colon Z \to X$ over $w$ ($p(\theta) = w$) such that $\varphi \circ \theta = \psi$:
\[
\forall Z\,\forall \psi\colon Z \to Y\,\forall w\colon p(Z) \to I\;\bigl(p(\psi) = u \circ w \implies \exists!\theta\colon Z \to X\;(p(\theta) = w \wedge \varphi \circ \theta = \psi)\bigr).
\]
Any two $p$-Cartesian lifts of $u\colon I \to J$ with codomain $Y \in \mathcal{E}_J$ are canonically isomorphic via a unique vertical isomorphism in $\mathcal{E}_I$. The functor $p\colon \mathcal{E} \to \mathcal{B}$ is a \emph{Grothendieck fibration} (or \emph{fibered category}) if every base morphism $u\colon I \to J$ and codomain object $Y \in \mathcal{E}_J$ admit a $p$-Cartesian lift $\operatorname{Cart}(u, Y)\colon u^*Y \to Y$; a chosen class-function assignment $(u, Y) \mapsto \operatorname{Cart}(u, Y)$ of such lifts is a \emph{cleavage}, which is \emph{split} when $1_J^* Y = Y$ and $(u \circ v)^* Y = v^*(u^* Y)$ hold on the nose. Given an indexed category (contravariant pseudofunctor) $\Phi\colon \mathcal{B}^{\mathrm{op}} \to \mathbf{Cat}$, the \emph{Grothendieck construction} builds the total category $\int_{\mathcal{B}} \Phi$ whose objects are pairs $\pair{I}{X}$ with $I \in \operatorname{Ob}(\mathcal{B})$ and $X \in \operatorname{Ob}(\Phi(I))$, and whose morphisms $\pair{u}{f}\colon \pair{I}{X} \to \pair{J}{Y}$ consist of $u\colon I \to J$ in $\mathcal{B}$ together with $f\colon X \to \Phi(u)(Y)$ in $\Phi(I)$, establishing the fundamental 2-categorical equivalence between cloven Grothendieck fibrations over $\mathcal{B}$ and indexed categories on $\mathcal{B}^{\mathrm{op}}$.

\subsection{Abelian Categories, Exact Sequences, and the Snake Lemma Connecting Homomorphism}
\label{subsec:abelian_categories_exact_sequences_snake_lemma}

An \emph{Abelian category} $\mathcal{A}$ (whose objects and morphism graphs are encoded via ordered pairs) is an additive category equipped with a zero object $0$, binary biproducts $A \oplus B$, kernels $\ker(f) = \operatorname{eq}(f, 0)$, and cokernels $\operatorname{coker}(f) = \operatorname{coeq}(f, 0)$ for all morphisms $f\colon A \to B$, such that every monomorphism is normal (the kernel of its cokernel), every epimorphism is conormal (the cokernel of its kernel), and for every morphism $f\colon A \to B$ the canonical comparison morphism from the coimage to the image is an isomorphism:
\[
\operatorname{coim}(f) = \operatorname{coker}(\ker(f)) \xrightarrow{\;\sim\;} \operatorname{im}(f) = \ker(\operatorname{coker}(f)).
\]
A composable pair $A \xrightarrow{f} B \xrightarrow{g} C$ is \emph{exact at $B$} when $\operatorname{im}(f) = \ker(g)$, and a sequence $0 \to A \xrightarrow{f} B \xrightarrow{g} C \to 0$ is a \emph{short exact sequence} when $f$ is a monomorphism ($\ker(f) = 0$), $g$ is an epimorphism ($\operatorname{coker}(g) = 0$), and $\operatorname{im}(f) = \ker(g)$. Given any commutative diagram of two short exact rows with vertical morphisms $(\alpha, \beta, \gamma)$, the \emph{Snake Lemma} constructs the natural connecting homomorphism $\delta\colon \ker(\gamma) \to \operatorname{coker}(\alpha)$ via the canonical zig-zag diagram chase $c \mapsto b \mapsto \beta(b) \mapsto a' \mapsto [a'] \in \operatorname{coker}(\alpha)$, yielding a six-term exact sequence
\[
\ker(\alpha) \longrightarrow \ker(\beta) \longrightarrow \ker(\gamma) \xrightarrow{\;\delta\;} \operatorname{coker}(\alpha) \longrightarrow \operatorname{coker}(\beta) \longrightarrow \operatorname{coker}(\gamma).
\]

\subsection{Chain Complexes, Projective and Injective Resolutions, and Derived Functors $\mathrm{Ext}^n$ and $\mathrm{Tor}_n$}
\label{subsec:chain_complexes_resolutions_derived_functors_ext_tor}

Let $\mathcal{A}$ be an Abelian category whose graded objects and differential graphs are represented via ordered pairs. A \emph{chain complex} $(C_\bullet, d_\bullet)$ in $\mathcal{A}$ consists of a $\mathbb{Z}$-indexed family of objects $C_n \in \operatorname{Ob}(\mathcal{A})$ and boundary differentials $d_n\colon C_n \to C_{n-1}$ satisfying the fundamental nilpotency identity $d_{n-1} \circ d_n = 0$ (equivalently $\operatorname{im}(d_{n+1}) \subseteq \ker(d_n)$), with $n$-th homology object $H_n(C_\bullet) = \ker(d_n)/\operatorname{im}(d_{n+1})$. Two chain maps $f_\bullet, g_\bullet\colon C_\bullet \to D_\bullet$ are \emph{chain homotopic} ($f_\bullet \simeq g_\bullet$) via a degree-$+1$ family $s_n\colon C_n \to D_{n+1}$ when $f_n - g_n = d^D_{n+1} \circ s_n + s_{n-1} \circ d^C_n$, inducing identical morphisms $H_n(f_\bullet) = H_n(g_\bullet)$ on homology. When $\mathcal{A}$ has enough projectives (respectively enough injectives), every object $A \in \operatorname{Ob}(\mathcal{A})$ admits a \emph{projective resolution} $P_\bullet \xrightarrow{\varepsilon} A \to 0$ (respectively an \emph{injective resolution} $0 \to A \xrightarrow{\eta} I^\bullet$) unique up to chain homotopy equivalence by the Comparison Theorem; applying the homological and cohomological functors yields the left derived torsion functors and right derived extension functors
\[
\operatorname{Tor}_n^{\mathcal{A}}(A, B) = H_n(P_\bullet \otimes B), \qquad
\operatorname{Ext}^n_{\mathcal{A}}(A, B) = H^n(\operatorname{Hom}_{\mathcal{A}}(P_\bullet, B)) \cong H^n(\operatorname{Hom}_{\mathcal{A}}(A, I^\bullet)),
\]
and every short exact sequence $0 \to A \to B \to C \to 0$ of chain complexes induces a natural \emph{long exact homology/cohomology sequence} via the Snake Lemma connecting homomorphisms $\partial_n\colon H_n(C_\bullet) \to H_{n-1}(A_\bullet)$.

\subsection{Triangulated Categories, Distinguished Triangles, the Octahedral Axiom, and Derived Categories}
\label{subsec:triangulated_categories_octahedral_derived_categories}

A \emph{triangulated category} $(\mathcal{D}, [1], \Delta)$ is an additive category $\mathcal{D}$ (whose objects, shift functors, and roof fractions are encoded via ordered pairs) equipped with an additive automorphism $[1]\colon \mathcal{D} \xrightarrow{\;\sim\;} \mathcal{D}$ (the suspension or translation functor $X \mapsto X[1]$) and a class $\Delta$ of \emph{distinguished triangles} $X \xrightarrow{u} Y \xrightarrow{v} Z \xrightarrow{w} X[1]$ satisfying the four Verdier axioms \textup{(TR1)}--\textup{(TR4)}:
\begin{enumerate}
\item \textup{(TR1)} $\Delta$ is closed under isomorphism, contains $X \xrightarrow{1_X} X \to 0 \to X[1]$ for every $X$, and embeds every morphism $u\colon X \to Y$ into a distinguished triangle via the mapping cone $\operatorname{Cone}(u)$;
\item \textup{(TR2)} \emph{Rotation}: $X \xrightarrow{u} Y \xrightarrow{v} Z \xrightarrow{w} X[1]$ is distinguished if and only if $Y \xrightarrow{v} Z \xrightarrow{w} X[1] \xrightarrow{-u[1]} Y[1]$ is distinguished;
\item \textup{(TR3)} \emph{Morphism extension}: any commutative square between the first two arrows of two distinguished triangles extends to a morphism of triangles;
\item \textup{(TR4)} \emph{Octahedral Axiom}: given composable morphisms $X \xrightarrow{u} Y \xrightarrow{v} Z$ and distinguished triangles on $u$, $v$, and $v \circ u$ with third objects $Z'$, $X'$, and $Y'$, there exists a distinguished triangle $Z' \to Y' \to X' \to Z'[1]$ making the octahedral diagram commute.
\end{enumerate}
Given an Abelian category $\mathcal{A}$ and its homotopy category of chain complexes $\mathbf{K}(\mathcal{A})$, the class $S$ of \emph{quasi-isomorphisms} (chain maps inducing isomorphisms on all homology objects $H_n$) forms a localizing multiplicative system in $\mathbf{K}(\mathcal{A})$ admitting a Gabriel--Zisman calculus of left and right fractions (where morphisms in the localization are represented by roofs $X_\bullet \xleftarrow{s} W_\bullet \xrightarrow{f} Y_\bullet$ with $s \in S$), yielding the \emph{derived category} $\mathbf{D}(\mathcal{A}) = \mathbf{K}(\mathcal{A})[S^{-1}]$ together with the universal localization functor $Q\colon \mathbf{K}(\mathcal{A}) \to \mathbf{D}(\mathcal{A})$.

\subsection{Quillen Model Categories, Weak Factorization Systems, and the Homotopy Category}
\label{subsec:quillen_model_categories_lifting_homotopy_category}

A \emph{Quillen model category} $(\mathcal{M}, \mathcal{W}, \mathcal{C}, \mathcal{F})$ is a complete and cocomplete category $\mathcal{M}$ (whose morphism classes, commutative lifting squares, and cylinder/path factorizations are encoded via ordered pairs) equipped with three distinguished classes of morphisms---weak equivalences $\mathcal{W}$, cofibrations $\mathcal{C}$, and fibrations $\mathcal{F}$ (with acyclic cofibrations $\mathcal{C} \cap \mathcal{W}$ and acyclic fibrations $\mathcal{F} \cap \mathcal{W}$)---satisfying the five Quillen axioms \textup{(MC1)}--\textup{(MC5)}:
\begin{enumerate}
\item \textup{(MC1)} \emph{Limits and colimits}: $\mathcal{M}$ has all finite limits and colimits (including initial object $\emptyset$ and terminal object $1$);
\item \textup{(MC2)} \emph{Two-out-of-three property}: for composable morphisms $f$ and $g$, if any two of $f$, $g$, and $g \circ f$ belong to $\mathcal{W}$, then so does the third;
\item \textup{(MC3)} \emph{Retracts}: each of $\mathcal{W}$, $\mathcal{C}$, and $\mathcal{F}$ is closed under retracts in the arrow category;
\item \textup{(MC4)} \emph{Lifting ($i \boxslash p$)}: for every commutative square with left vertical arrow $i\colon A \to B$ in $\mathcal{C}$ and right vertical arrow $p\colon X \to Y$ in $\mathcal{F}$, a diagonal lift $h\colon B \to X$ satisfying $h \circ i = f$ and $p \circ h = g$ exists whenever either $i \in \mathcal{C} \cap \mathcal{W}$ or $p \in \mathcal{F} \cap \mathcal{W}$;
\item \textup{(MC5)} \emph{Weak factorization systems}: every morphism $f\colon X \to Y$ factors both as $f = p \circ i$ with $i \in \mathcal{C}$ and $p \in \mathcal{F} \cap \mathcal{W}$, and as $f = q \circ j$ with $j \in \mathcal{C} \cap \mathcal{W}$ and $q \in \mathcal{F}$.
\end{enumerate}
An object $X \in \operatorname{Ob}(\mathcal{M})$ is \emph{cofibrant} if $\emptyset \to X$ is in $\mathcal{C}$, \emph{fibrant} if $X \to 1$ is in $\mathcal{F}$, and \emph{bifibrant} if both hold. By Quillen's Whitehead Theorem, the localization $\operatorname{Ho}(\mathcal{M}) = \mathcal{M}[\mathcal{W}^{-1}]$ is canonically equivalent to the quotient of the full subcategory $\mathcal{M}_{cf}$ of bifibrant objects modulo cylinder/path homotopy $\sim$, and every \emph{Quillen adjunction} $L \dashv R\colon \mathcal{M} \rightleftarrows \mathcal{N}$ (where $L$ preserves cofibrations and acyclic cofibrations, equivalently $R$ preserves fibrations and acyclic fibrations) induces a total derived adjunction $\mathbf{L}L \dashv \mathbf{R}R\colon \operatorname{Ho}(\mathcal{M}) \rightleftarrows \operatorname{Ho}(\mathcal{N})$ which is an equivalence of homotopy categories when $(L, R)$ is a Quillen equivalence.

\subsection{Simplicial Sets, Kan Complexes, Quasi-Categories, and the Homotopy Coherent Nerve}
\label{subsec:simplicial_sets_kan_complexes_quasi_categories}

Let $\boldsymbol{\Delta}$ denote the simplex category of non-empty finite linearly ordered ordinals $[n] = \{0 < 1 < \dots < n\}$ ($n \ge 0$) with order-preserving maps generated by coface maps $\delta^i\colon [n-1] \to [n]$ and codegeneracy maps $\sigma^j\colon [n+1] \to [n]$. A \emph{simplicial set} $X_\bullet\colon \boldsymbol{\Delta}^{\mathrm{op}} \to \mathbf{Set}$ (whose graded $n$-simplices $X_n$ and face/degeneracy operator graphs $d_i\colon X_n \to X_{n-1}$, $s_j\colon X_n \to X_{n+1}$ are encoded via ordered pairs) satisfies the standard \emph{simplicial identities}:
\[
d_i \circ d_j = d_{j-1} \circ d_i \quad (i < j), \qquad
s_i \circ s_j = s_{j+1} \circ s_i \quad (i \le j), \qquad
d_i \circ s_j = \begin{cases}
s_{j-1} \circ d_i & \text{if } i < j, \\
\mathrm{id} & \text{if } i = j \text{ or } i = j+1, \\
s_j \circ d_{i-1} & \text{if } i > j+1.
\end{cases}
\]
For each $n \ge 1$ and $0 \le k \le n$, the \emph{$k$-th horn} $\Lambda^n_k \subseteq \Delta^n$ is the sub-simplicial set of the standard $n$-simplex $\Delta^n = \boldsymbol{\Delta}(-, [n])$ obtained by removing the interior of $\Delta^n$ and the $k$-th face $d_k(\Delta^n)$; the horn $\Lambda^n_k$ is an \emph{inner horn} when $0 < k < n$ and an \emph{outer horn} when $k = 0$ or $k = n$. A simplicial set $X_\bullet$ is:
\begin{enumerate}
\item a \emph{quasi-category} (or \emph{$\infty$-category} in the Joyal--Boardman--Vogt model) if every inner horn $\Lambda^n_k \to X_\bullet$ ($0 < k < n$) admits a filler $\Delta^n \to X_\bullet$;
\item a \emph{Kan complex} (or \emph{$\infty$-groupoid}, the fibrant objects of the Quillen--Kan model structure) if every horn $\Lambda^n_k \to X_\bullet$ ($0 \le k \le n$, including outer horns) admits a filler $\Delta^n \to X_\bullet$, so every Kan complex is a fortiori a quasi-category;
\item the \emph{categorical nerve} $\operatorname{N}(\mathcal{C})_\bullet$ of a small 1-category $\mathcal{C}$ (whose $n$-simplices are composable $n$-tuples of morphisms $C_0 \to C_1 \to \dots \to C_n$) if and only if every inner horn $\Lambda^n_k \to \operatorname{N}(\mathcal{C})_\bullet$ ($0 < k < n$) admits a \emph{unique} filler, while Cordier's \emph{homotopy coherent nerve} $\operatorname{N}_{\mathrm{hc}}(\mathcal{C}_\Delta)$ extends this construction to simplicially enriched categories $\mathcal{C}_\Delta$ as the right Quillen adjoint to rigidification $\mathfrak{C} \dashv \operatorname{N}_{\mathrm{hc}}$.
\end{enumerate}

\subsection{Grothendieck Sites, Covering Sieves, Sheaf Gluing, and Plus-Construction Sheafification}
\label{subsec:grothendieck_sites_sieves_sheaf_gluing_sheafification}

Let $\mathcal{C}$ be a category whose objects, morphism graphs, and sieve subfunctors are encoded via ordered pairs. For an object $U \in \operatorname{Ob}(\mathcal{C})$, a \emph{sieve on $U$} is a right ideal of morphisms with codomain $U$ (equivalently a subfunctor $S \hookrightarrow h_U = \operatorname{Hom}_{\mathcal{C}}(-, U)$ closed under precomposition: $f \in S(V) \implies f \circ g \in S(W)$ for all $g\colon W \to V$). Given a morphism $h\colon V \to U$, the \emph{pullback sieve} $h^*(S)$ on $V$ is defined by $h^*(S)(W) = \{g\colon W \to V \mid h \circ g \in S(W)\}$. A \emph{Grothendieck topology} $J$ on $\mathcal{C}$ assigns to each object $U$ a collection $J(U)$ of \emph{covering sieves} on $U$ satisfying:
\begin{enumerate}
\item \emph{Maximal sieve axiom}: the maximal sieve $t_U = \operatorname{Hom}_{\mathcal{C}}(-, U)$ belongs to $J(U)$;
\item \emph{Pullback stability}: if $S \in J(U)$ and $h\colon V \to U$, then $h^*(S) \in J(V)$;
\item \emph{Local character (transitivity)}: if $S \in J(U)$ and $R$ is any sieve on $U$ such that $h^*(R) \in J(V)$ for all $h \in S(V)$, then $R \in J(U)$.
\end{enumerate}
The pair $(\mathcal{C}, J)$ is called a \emph{Grothendieck site}. Given a presheaf $P\colon \mathcal{C}^{\mathrm{op}} \to \mathbf{Set}$ and a sieve $S \in J(U)$, a \emph{matching family} for $P$ on $S$ is a natural transformation $x\colon S \Rightarrow P$ (a family $\{x_f \in P(\operatorname{dom}(f))\}_{f \in S}$ satisfying $P(g)(x_f) = x_{f \circ g}$), and an \emph{amalgamation} is an element $s \in P(U)$ such that $P(f)(s) = x_f$ for all $f \in S$. The presheaf $P$ is a \emph{$J$-separated presheaf} if every matching family has at most one amalgamation, and a \emph{$J$-sheaf} if every matching family on every covering sieve $S \in J(U)$ has a \emph{unique} amalgamation ($\operatorname{Hom}(h_U, P) \xrightarrow{\sim} \operatorname{Hom}(S, P)$). Grothendieck's \emph{plus-construction} $P^+(U) = \varinjlim_{S \in J(U)^{\mathrm{op}}} \operatorname{Match}(S, P)$ maps any presheaf $P$ to a separated presheaf $P^+$ and maps any separated presheaf to a sheaf, so the iterated functor $\mathbf{a}(P) = P^{++}\colon \mathbf{PSh}(\mathcal{C}) \to \mathbf{Sh}(\mathcal{C}, J)$ is the exact left adjoint (\emph{sheafification}) to the full inclusion $\mathbf{Sh}(\mathcal{C}, J) \hookrightarrow \mathbf{PSh}(\mathcal{C})$, yielding the Grothendieck topos $\mathbf{Sh}(\mathcal{C}, J)$.

\subsection{Spectral Sequences, Filtered Complexes, Exact Couples, and Grothendieck Convergence}
\label{subsec:spectral_sequences_filtered_complexes_grothendieck}

Let $\mathcal{A}$ be an Abelian category whose bigraded objects, differentials, and filtration stages are encoded via ordered pairs. A \emph{cohomological spectral sequence} in $\mathcal{A}$ (starting at page $r_0 \ge 0$) consists of a family of bigraded objects $\{E_r^{p,q}\}_{p,q \in \mathbb{Z},\, r \ge r_0}$ together with differentials $d_r^{p,q}\colon E_r^{p,q} \to E_r^{p+r, q-r+1}$ of bidegree $(r, 1-r)$ satisfying $d_r \circ d_r = 0$, together with canonical page-transition isomorphisms
\[
E_{r+1}^{p,q} \;\cong\; H^{p,q}(E_r^{\bullet,\bullet}, d_r) \;=\; \frac{\ker\bigl(d_r^{p,q}\colon E_r^{p,q} \to E_r^{p+r, q-r+1}\bigr)}{\operatorname{im}\bigl(d_r^{p-r, q+r-1}\colon E_r^{p-r, q+r-1} \to E_r^{p,q}\bigr)}.
\]
Equivalently, any \emph{exact couple} $(D, E, i, j, k)$ in $\mathcal{A}$---an exact triangle $D \xrightarrow{i} D \xrightarrow{j} E \xrightarrow{k} D$---generates a derived exact couple $(D', E', i', j', k')$ with $D' = \operatorname{im}(i)$, $E' = H(E, j \circ k)$, and differential $d = j \circ k\colon E \to E$ satisfying $d \circ d = j \circ (k \circ j) \circ k = 0$. Given a bounded decreasingly \emph{filtered cochain complex} $\cdots \supseteq F^p C^\bullet \supseteq F^{p+1} C^\bullet \supseteq \cdots$ (exhaustive and Hausdorff), the spectral sequence with $E_0^{p,q} = \operatorname{Gr}^p C^{p+q} = F^p C^{p+q}/F^{p+1} C^{p+q}$ and $E_1^{p,q} = H^{p+q}(\operatorname{Gr}^p C^\bullet)$ converges strongly to the filtered cohomology abutment:
\[
E_1^{p,q} \;\Longrightarrow\; H^{p+q}(C^\bullet), \qquad E_\infty^{p,q} \;\cong\; \operatorname{Gr}^p H^{p+q}(C^\bullet) = \frac{F^p H^{p+q}(C^\bullet)}{F^{p+1} H^{p+q}(C^\bullet)}.
\]
In particular, if $F\colon \mathcal{A} \to \mathcal{B}$ and $G\colon \mathcal{B} \to \mathcal{C}$ are left exact functors between Abelian categories with enough injectives such that $F$ sends injective objects of $\mathcal{A}$ to $G$-acyclic objects of $\mathcal{B}$ ($(\mathbf{R}^q G)(F(I)) = 0$ for all $q > 0$), the \emph{Grothendieck composite-functor spectral sequence} satisfies
\[
E_2^{p,q} = (\mathbf{R}^p G)\bigl((\mathbf{R}^q F)(A)\bigr) \;\Longrightarrow\; \mathbf{R}^{p+q}(G \circ F)(A).
\]

\subsection{Symmetric Monoidal Closed Categories, Enriched Yoneda, and Day Convolution}
\label{subsec:symmetric_monoidal_closed_enriched_yoneda_day_convolution}

Let $\mathcal{V} = (\mathcal{V}_0, \otimes, I, \alpha, \lambda, \rho, \sigma)$ be a \emph{symmetric monoidal category} whose objects, tensor product graphs, coherence natural isomorphisms, and enriched hom-objects are encoded via ordered pairs. By \emph{Mac Lane's Coherence Theorem}, every formal diagram built from the associator $\alpha_{A,B,C}\colon (A \otimes B) \otimes C \xrightarrow{\sim} A \otimes (B \otimes C)$, left and right unitors $\lambda_A\colon I \otimes A \xrightarrow{\sim} A$ and $\rho_A\colon A \otimes I \xrightarrow{\sim} A$, and involutive braiding $\sigma_{A,B}\colon A \otimes B \xrightarrow{\sim} B \otimes A$ ($\sigma_{B,A} \circ \sigma_{A,B} = \operatorname{id}_{A \otimes B}$) commutes whenever the \emph{pentagon identity}
\[
(\operatorname{id}_A \otimes \alpha_{B,C,D}) \circ \alpha_{A,B \otimes C,D} \circ (\alpha_{A,B,C} \otimes \operatorname{id}_D) \;=\; \alpha_{A,B,C \otimes D} \circ \alpha_{A \otimes B,C,D}
\]
and the \emph{triangle identity} $(\operatorname{id}_A \otimes \lambda_B) \circ \alpha_{A,I,B} = \rho_A \otimes \operatorname{id}_B$ (together with hexagon coherence) hold. Furthermore, $\mathcal{V}$ is \emph{symmetric monoidal closed} if for every object $B \in \operatorname{Ob}(\mathcal{V})$, the endofunctor $(-) \otimes B\colon \mathcal{V} \to \mathcal{V}$ admits a right adjoint $[B, -]\colon \mathcal{V} \to \mathcal{V}$ (\emph{internal hom}), yielding natural tensor-hom adjunction isomorphisms
\[
\operatorname{Hom}_{\mathcal{V}}(A \otimes B, C) \;\cong\; \operatorname{Hom}_{\mathcal{V}}\bigl(A, [B, C]\bigr), \qquad [A \otimes B, C] \;\cong\; \bigl[A, [B, C]\bigr].
\]
When $\mathcal{V}$ is complete and cocomplete symmetric monoidal closed, a \emph{$\mathcal{V}$-enriched category} $\underline{\mathcal{C}}$ consists of a class of objects $\operatorname{Ob}(\underline{\mathcal{C}})$, hom-objects $\underline{\mathcal{C}}(X, Y) \in \operatorname{Ob}(\mathcal{V})$, composition morphisms $M_{X,Y,Z}\colon \underline{\mathcal{C}}(Y, Z) \otimes \underline{\mathcal{C}}(X, Y) \to \underline{\mathcal{C}}(X, Z)$, and unit elements $j_X\colon I \to \underline{\mathcal{C}}(X, X)$ satisfying enriched associativity and unit laws. For any $\mathcal{V}$-functor $F\colon \underline{\mathcal{C}} \to \mathcal{V}$, the \emph{enriched Yoneda Lemma} establishes an isomorphism in $\mathcal{V}$ via the enriched end
\[
F(A) \;\cong\; [\underline{\mathcal{C}}, \mathcal{V}]\bigl(\underline{\mathcal{C}}(A, -), F\bigr) \;=\; \int_{X \in \underline{\mathcal{C}}} \bigl[\underline{\mathcal{C}}(A, X), F(X)\bigr].
\]
Finally, if $\underline{\mathcal{C}}$ is a small symmetric monoidal $\mathcal{V}$-category, the enriched functor category $[\underline{\mathcal{C}}^{\mathrm{op}}, \mathcal{V}]$ inherits a symmetric monoidal closed structure via \emph{Day convolution} coend
\[
(F \star_{\mathrm{Day}} G)(Z) \;=\; \int^{X, Y \in \underline{\mathcal{C}}} \underline{\mathcal{C}}(Z, X \otimes_{\mathcal{C}} Y) \otimes F(X) \otimes G(Y),
\]
such that the enriched Yoneda embedding $\mathbf{y}\colon \underline{\mathcal{C}} \hookrightarrow [\underline{\mathcal{C}}^{\mathrm{op}}, \mathcal{V}]$ is strong symmetric monoidal: $\mathbf{y}(A \otimes_{\mathcal{C}} B) \cong \mathbf{y}(A) \star_{\mathrm{Day}} \mathbf{y}(B)$.

\subsection{Ends, Coends, the Fubini Interchange Theorem, and Weighted Limits}
\label{subsec:ends_coends_fubini_weighted_limits}

Let $\mathcal{C}$ and $\mathcal{D}$ be categories whose objects, morphisms, and bifunctor graphs $T\colon \mathcal{C}^{\mathrm{op}} \times \mathcal{C} \to \mathcal{D}$ are encoded via ordered pairs. An \emph{extranatural wedge} from an object $W \in \operatorname{Ob}(\mathcal{D})$ to $T$ is a family of morphisms $\omega = (\omega_X\colon W \to T(X, X))_{X \in \operatorname{Ob}(\mathcal{C})}$ satisfying the dinaturality condition $T(\operatorname{id}_X, f) \circ \omega_X = T(f, \operatorname{id}_Y) \circ \omega_Y$ in $T(X, Y)$ for every morphism $f\colon X \to Y$ in $\mathcal{C}$. An \emph{end} of $T$, denoted $\int_{X \in \mathcal{C}} T(X, X)$, is a terminal wedge $(E, \pi_X\colon E \to T(X, X))$ characterized by the equalizer diagram (when $\mathcal{C}$ is small and $\mathcal{D}$ is complete)
\[
\int_{X \in \mathcal{C}} T(X, X) \;\hookrightarrow\; \prod_{X \in \operatorname{Ob}(\mathcal{C})} T(X, X) \;\rightrightarrows\; \prod_{f\colon X \to Y} T(X, Y).
\]
Dually, a \emph{coend} $\int^{X \in \mathcal{C}} T(X, X)$ is an initial cowedge $(\iota_X\colon T(X, X) \to C, C)$ characterized as the coequalizer of $\coprod_{f\colon X \to Y} T(Y, X) \rightrightarrows \coprod_{X \in \operatorname{Ob}(\mathcal{C})} T(X, X)$. In particular, the class of natural transformations between functors $F, G\colon \mathcal{C} \to \mathcal{D}$ is canonically represented as the end $\operatorname{Nat}(F, G) \cong \int_{X \in \mathcal{C}} \operatorname{Hom}_{\mathcal{D}}(F(X), G(X))$, and composition of profunctors (distributors) $P\colon \mathcal{A}^{\mathrm{op}} \times \mathcal{B} \to \mathbf{Set}$ and $Q\colon \mathcal{B}^{\mathrm{op}} \times \mathcal{C} \to \mathbf{Set}$ is given by the coend $(Q \odot P)(A, C) = \int^{B \in \mathcal{B}} P(A, B) \times Q(B, C)$. By the \emph{Fubini Interchange Theorem} for iterated ends and coends of $S\colon (\mathcal{A} \times \mathcal{B})^{\mathrm{op}} \times (\mathcal{A} \times \mathcal{B}) \to \mathcal{D}$, when the inner ends (respectively coends) exist, there are canonical isomorphisms
\[
\int_{A \in \mathcal{A}} \int_{B \in \mathcal{B}} S(A, B, A, B) \;\cong\; \int_{(A, B) \in \mathcal{A} \times \mathcal{B}} S(A, B, A, B) \;\cong\; \int_{B \in \mathcal{B}} \int_{A \in \mathcal{A}} S(A, B, A, B),
\]
and dually for iterated coends $\int^{A} \int^{B} \cong \int^{(A,B)} \cong \int^{B} \int^{A}$. Finally, in a $\mathcal{V}$-enriched category $\underline{\mathcal{D}}$ over a symmetric monoidal closed category $\mathcal{V}$, the \emph{$W$-weighted limit} $\{W, F\}$ of a $\mathcal{V}$-diagram $F\colon \underline{\mathcal{J}} \to \underline{\mathcal{D}}$ weighted by $W\colon \underline{\mathcal{J}} \to \mathcal{V}$ and the \emph{$W'$-weighted colimit} $W' \star F$ (with $W'\colon \underline{\mathcal{J}}^{\mathrm{op}} \to \mathcal{V}$) are represented by the enriched ends
\[
\underline{\mathcal{D}}\bigl(D, \{W, F\}\bigr) \;\cong\; \int_{J \in \underline{\mathcal{J}}} \bigl[W(J),\, \underline{\mathcal{D}}(D, F(J))\bigr], \qquad \underline{\mathcal{D}}\bigl(W' \star F, D\bigr) \;\cong\; \int_{J \in \underline{\mathcal{J}}} \bigl[W'(J),\, \underline{\mathcal{D}}(F(J), D)\bigr].
\]

\subsection{Bicategories, 2-Cell Interchange, Eckmann-Hilton Commutativity, and B\'enabou Coherence}
\label{subsec:bicategories_interchange_eckmann_hilton_benabou}

A \emph{bicategory} $\mathcal{B}$ (introduced by B\'enabou) whose $0$-cells (objects), $1$-cells (morphisms), and $2$-cells (globular transformations $\alpha\colon f \Rightarrow g$) are encoded via ordered pairs consists of a class of $0$-cells $\operatorname{Ob}(\mathcal{B})$, hom-categories $\mathcal{B}(A, B)$ whose objects are $1$-cells $f\colon A \to B$ and whose morphisms are $2$-cells $\alpha\colon f \Rightarrow g$ under strict \emph{vertical composition} $\beta \cdot \alpha$, horizontal composition functors $c_{A,B,C}\colon \mathcal{B}(B, C) \times \mathcal{B}(A, B) \to \mathcal{B}(A, C)$ written $(g, f) \mapsto g \circ f$ and $(\beta, \alpha) \mapsto \beta \ast \alpha$, identity $1$-cells $\operatorname{id}_A \in \mathcal{B}(A, A)$, and natural $2$-cell isomorphisms $\alpha_{h,g,f}\colon (h \circ g) \circ f \xRightarrow{\sim} h \circ (g \circ f)$, $\lambda_f\colon \operatorname{id}_B \circ f \xRightarrow{\sim} f$, and $\rho_f\colon f \circ \operatorname{id}_A \xRightarrow{\sim} f$ satisfying the pentagon and triangle coherence axioms. Functoriality of horizontal composition is equivalent to the \emph{middle-four interchange law} for composable $2$-cells $\alpha\colon f_1 \Rightarrow f_2$, $\alpha'\colon f_2 \Rightarrow f_3$ in $\mathcal{B}(A, B)$ and $\beta\colon g_1 \Rightarrow g_2$, $\beta'\colon g_2 \Rightarrow g_3$ in $\mathcal{B}(B, C)$:
\[
(\beta' \cdot \beta) \ast (\alpha' \cdot \alpha) \;=\; (\beta' \ast \alpha') \cdot (\beta \ast \alpha).
\]
By the \emph{Eckmann-Hilton argument}, horizontal and vertical composition on the endomorphism $2$-cell monoid $\operatorname{End}_{\mathcal{B}(A, A)}(\operatorname{id}_A) = \mathcal{B}(A, A)(\operatorname{id}_A, \operatorname{id}_A)$ coincide and are commutative: $\beta \ast \alpha = \beta \cdot \alpha = \alpha \cdot \beta = \alpha \ast \beta$. An \emph{internal adjunction} $f \dashv u\colon A \leftrightarrows B$ in $\mathcal{B}$ consists of $1$-cells $f\colon A \to B$ and $u\colon B \to A$ together with unit $\eta\colon \operatorname{id}_A \Rightarrow u \circ f$ and counit $\varepsilon\colon f \circ u \Rightarrow \operatorname{id}_B$ satisfying the triangle identities $(\varepsilon \ast \operatorname{id}_f) \cdot \alpha_{f,u,f}^{-1} \cdot (\operatorname{id}_f \ast \eta) = \operatorname{id}_f$ and $(\operatorname{id}_u \ast \varepsilon) \cdot \alpha_{u,f,u} \cdot (\eta \ast \operatorname{id}_u) = \operatorname{id}_u$ (up to unitors). Finally, \emph{B\'enabou's Coherence and Strictification Theorem} establishes that every formal diagram of canonical constraint $2$-cells commutes and every bicategory $\mathcal{B}$ is biequivalent (via the Yoneda pseudofunctor $\mathcal{B} \to \mathbf{Cat}(\mathcal{B}^{\mathrm{op}}, \mathbf{Cat})$) to a strict $2$-category $\operatorname{Strict}(\mathcal{B})$.

\ifdefined\qwertyflag
\undef\qwertyflag
\end{document}
\fi
